\documentclass[conference]{IEEEtran}
\IEEEoverridecommandlockouts
% The preceding line is only needed to identify funding in the first footnote. If that is unneeded, please comment it out.
\usepackage{cite}
\usepackage{amsmath,amssymb,amsfonts}
\usepackage{algorithmic}
\usepackage{graphicx}
\usepackage{textcomp}
\usepackage{xcolor}
\usepackage[skip=1em, indent=1em]{parskip}
\def\BibTeX{{\rm B\kern-.05em{\sc i\kern-.025em b}\kern-.08em
    T\kern-.1667em\lower.7ex\hbox{E}\kern-.125emX}}
\begin{document}

\title{How Social Media Algorithms Manipulate Users and Push Harmful Content\\}

\author{\IEEEauthorblockN{Marlen Vining}}

\maketitle 

\begin{abstract}

Social media platforms primarily generate revenue from the collection of user data and targeted advertising. This is why it is in their best interest to keep users engaged as long as possible. This is done by creating detailed profiles on users to find what specific kinds of content appeal to them most. This is done by collecting data on what posts users are most interested in regardless of if the interest is positive or not. Many algorithms are designed to serve content to users that will keep them using the platform. Many of these algorithms purposefully serve content to users that is designed to trigger strong emotional responses. This is done to extend the time users spend on the platform and create a dependence. This often has a direct negative impact on the mental health and well being of users. The goal of this paper is to understand the tactics used by these platforms and uncover the harm it causes. It will examine how these algorithms decide what content to serve to users, as well as go over cases where algorithms have potentially caused real harm to users. 

\end{abstract}

\section{How Social Media Algorithms Work}

Most major social media platforms receive millions of posts per hour from billions of active users. "The reason why social media platforms use algorithms is to more organically filter through the large amount of content that is available on each platform. Algorithms do the work of delivering content that is potentially more “interesting” for a user to the detriment of posts which are deemed irrelevant or low-quality - either in general, or to a specific user."\cite{algorithms2021} Not only would it be impossible for anyone to keep up with every post made on even one of these sites, the majority of this content would be both uninteresting to, and not made for many of the users on the platform. This is why platforms have implemented algorithms to show only the posts that users are likely to engage with. This is done by collecting data on user behavior to create a curated feed for each individual user based on what keeps them interested. As users spend more time on the platform, a more detailed profile is built on them in order to tune the algorithm.

These algorithms use a variety of factors to determine what content to show users. The specific metrics and the way they are prioritized differ from site to site, they generally include direct user actions such as likes and comments, and passive signals such as how much time a user spends on a specific post. Content is ranked by both what has a wide appeal, and what will appeal to the specific user viewing the platform. Data from data brokers also allows for platforms to create detailed profiles on user activity from a variety of sites. This method of curating content heavily prioritizes engagement over truth and quality in order to keep users on the platform as long as possible.

Modern social media algorithms are based on machine learning models. These take in data on users and make predictions on what content they will find most interesting. Content is classified by factors such as the target demographic and appeal of the post. When a user views their feed, they are given an up-to-date list of posts that align with what they typically engage with. Their reaction to this content is then measured and fed back into the algorithm to adjust future recommendations. Most platforms prioritize content that is both new, and popular. This provides users with an endless feed of content that is typically interesting and of decent quality. Intentional biases may also be set by the platform to limit the reach of content that violates terms of service as well as promote content that aligns with the purpose of the platform. These algorithms are constantly adjusted to be more effective and meet the needs of the platform. These models are often able to predict the likelihood that a user will interact with a post before they even see it.

\section{Why Algorithms Push Harmful Content}
Social media sites have a financial incentive to keep users on the platform as long as possible. These sites are also incentivized to push content that increases the effectiveness of advertising on the platform. "Algorithms  designed  to  enhance  user  engagement  have  significantly  shaped  how content  is  prioritized, recommended, and disseminated. While these algorithms aim  to  create a  personalized experience, they often operate  without  sufficient safeguards against the spread of  harmful content,  including racial  and ethnic misinformation."\cite{Lucas2024} Content that produces a strong emotional response in the user, as well as content that builds a dependence on the platform keeps users engaged longer and creates a pattern of repeated use. Because of this, algorithms are often designed to push these kinds of contents. While this may be an unintentional result, it can still cause real harm.

Negative emotional responses are often stronger and more impactful than positive ones. This leads controversial or emotionally damaging. People often feel influenced to spend more time on and share posts that cause a strong negative response. This leads to the spread of misinformation, conspiracy theories, hateful content, and other kinds of harmful posts. This not only means that this kind of content can rise in the algorithm faster than other kinds of posts, but also means that users that are drawn in by it may be recommended similar kinds of content. When someone gets pulled into a cycle of negative content, they are not only increasing the chances that they will continue to see this kind of content, but unknowingly contributing to the engagement metrics of these posts. This causes them to spread further and impact more people. Negative reactions drive discussion, which often leads to the most controversial statements from both sides being taken and reposted, continuing the cycle. On platforms where posts can be monetized, it is often in the interest of both the platform and the users to cause controversy for attention.

While most platforms ban the posting of things such as gore and extremist content, moderation is often lacking and easy for bad actors to bypass. In some cases this content is even altered to appear child friendly. One notable incident is when gore content was spread on platforms such as Instagram and Tik Tok using an AI filter to make the videos appear like animated clips from the Minions movies.\cite{Maiberg2025_Minions} These videos would often lead children to the source of the videos and expose them to unfiltered gore. Because of how the videos were filtered, it was not detected by moderation until the public discovered the trend and exposed it.

\section{Echo Chambers and Polarization}
Even posts that do not cause negative emotional responses can be harmful. In order to strengthen a reliance on the platform, many social media sites create an unrealistic level of reinforcement to views the user may have. This creates a strong sense of community and can lead to users feeling rejected by society. This leads users into "echo chambers" where they primarily see content that aligns with their views even if those views are only held by a minority of the total population. Posts from opposing views are typically framed as something to be criticized or ridiculed. "The echo chamber effect is strongest on social media sites that rely on algorithms to feed user content. Social media users tend to prefer information that conforms to their preexisting viewpoints and ignore dissenting information."\cite{Ktepi2021_EchoChamber}

Echo chambers lead to polarization. "When a larger community—such as a nation—consists of especially strictly delineated tribal groups, the result is increasing polarized political discourse, in which common ground becomes more difficult to find."\cite{Ktepi2021_EchoChamber} Because the viewpoints of the group are viewed as objective truth and rarely if ever challenged, anyone outside the group is viewed as objectively wrong or as an enemy. This leads to further conflict and increased engagement, causing people on all sides to form stronger dependencies on their echo chambers and withdraw further from opposing groups. Users that participate in echo chambers long enough start to believe that their views are held by the majority. This builds unrealistic expectations for the real world. It also leads to these users avoiding interactions with people of opposing groups both online and in person.

Another negative effect of echo chambers is that they often lead to a "pipeline" effect. With the lack of opposing views, more extreme versions of a group's views are more likely to be accepted. These extreme groups often make themselves appealing users who are already polarized enough to feel rejected by wider society. They provide a strong sense of community while conditioning members to believe more extreme views. This type of radicalization has been made extremely effective by social media, as these groups can reach vulnerable people through the internet. These groups contribute heavily to the spread of misinformation and benefit from the lack of pushback found in echo chambers.

\section{Misinformation and Disinformation}
Misinformation and disinformation have become a major issue in recent years. It has been used to influence politics and can be easily found on all social media platforms. Misinformation refers to any false information being spread. Disinformation is when that content is spread intentionally to convince people of a lie. While most misinformation on the internet today started out as disinformation, its nature allows for it to be spread by unknowing users. This increases its perceived legitimacy. When this content is shared or re-posted by people a user trusts, or becomes popular enough to be seen multiple times, it becomes easier to believe and continue spreading. Algorithms spread content based on engagement, not truth. "The flow of misinformation on Twitter is thus a function of both human and technical factors. Human biases play an important role: Since we’re more likely to react to content that taps into our existing grievances and beliefs, inflammatory tweets will generate quick engagement. It’s only after that engagement happens that the technical side kicks in: If a tweet is retweeted, favorited, or replied to by enough of its first viewers, the newsfeed algorithm will show it to more users, at which point it will tap into the biases of those users too—prompting even more engagement, and so on."\cite{Meserole2018_MisinformationSpread}

Misinformation can be used to alter or strengthen the beliefs of the users who view it. By presenting false information as true, it can strengthen the overall argument the information is meant to support. Misinformation has been used widely in political campaigns to make one side seem significantly better or worse than the other. People who get caught up in misinformation may form views that are difficult to alter even after the misinformation has been widely disproven. This is why it is such an effective tool for bad actors looking to influence people. Botting can be used to boost new information in the algorithm. This is when fake accounts are used to add likes, comments, and views to a post. This makes the post look as if it has real interaction. Not only does the algorithm see botted posts as popular, but the fake engagement can make them appear more legitimate to users, especially when comments are posted supporting it.

\section{The Impact On Users}
The way social media algorithms work can lead to real world harm if vulnerable users are shown the wrong kind of content. Content recommendations are based on the specific content that will keep a user engaged the longest. Vulnerable groups such as young people or people struggling with mental health issues may be more easily drawn in by content that could negatively impact them. This can lead to a negative spiral as engagement with these types of posts leads to more of them being recommended, flooding the user's feed with similar content. There have been many documented incidents of widespread harm caused by social media algorithms.

In a study done by Meta to measure the impacts of certain content on teens, "the 223 teens who often felt bad about their bodies after viewing Instagram, “eating disorder adjacent content” made up 10.5\% of what they saw on the platform. Among the other teens in the study, such content accounted for only 3.3\% of what they saw."\cite{Reuters2025_InstagramEDContent} Content involving harmful behavior such as self harm or eating disorders can trigger an increase in that behavior in people who are currently struggling with it or have in the past. "A scoping review from 2021 concludes 4 major factors of influence that eventually lead to disordered eating behaviors (13). These include visual appeal, content dissemination, socialized digital connections and adolescent marketer influencers (13)."\cite{Nawaz2024_SocialMediaED} Pushing this kind of content to users who are vulnerable to it can lead to direct harm of those users and the spread of that content to others who may be impacted by it.

When a French court investigated the platform Tik Tok over concerns that it was having a negative psychological effect on youth, it was noted that, "insufficient moderation of TikTok, its ease of access by minors and its sophisticated algorithm, which could push vulnerable individuals toward suicide by quickly trapping them in a loop of dedicated content."\cite{Reuters2025_TikTokSuicide} This illustrates how harmful and dangerous content can take over a user's feed. Interacting with one post, even out of a morbid curiosity, can be seen by the algorithm as a legitimate interest. This leads to similar content being recommended, trapping users in a cycle that can be hard to get out of. Being surrounded with this content can normalize negative thoughts and lead to users doing real harm to themselves or others.

Another more subtle way that algorithms cause harm is through addiction. Because engagement is the primary goal of a social media platform, they are designed to create a cycle of addiction to keep users coming back to the platform. With a variety of platforms competing for the time and data of their users, each platform needs to be as interesting and addictive as possible to keep the attention of their users. "Variably rewarding users with stimuli (likes, notifications, comments, etc.) keeps them engaged with content. When a user’s photo receives a “like,” the same dopamine pathways involved in motivation, reward, and addiction are activated."\cite{Qiu2021_PsychiatristsAlgorithms}

There are many negative effects to social media addiction that may be subtle enough to not immediately notice. "Teenagers who use social media regularly are more inclined to struggle with anxiety, depression, and a general decline in life satisfaction, according to multiple research studies. The long-term use of social media could negatively impact sleep quality, which exacerbates mental health problems, according to studies [7,15,16]. Adolescent eating disorders and appearance-related anxiety may worsen as an effect of the emphasis on ideals of beauty that permeate social media culture [7,17,18]."\cite{Amirthalingam2024_SocialMediaAddiction} Social media addiction can have a serious impact on a person's life. However, it is often the goal of these platforms to keep users addicted. Users that regularly spend long periods of time on the platform generate more revenue for the site.

\section{Targeted Advertising}
The primary way social media platforms generate revenue is through the use of targeted advertisements. This is when platforms use the profile they have constructed on a user to show them advertisements for products they are likely to be interested in. The process for this is similar to how normal posts are shown to users. Platforms make more money if users view more ads and if they can make the ads more effective. Targeted advertisements are designed to manipulate users and are often designed to target specific groups based on the detailed profiles built on users. Everything from age and location to gender and political alignment are collected to make advertising as effective as possible. This level of data collection is a major privacy concern.

One notable instance of this is when Meta was found to be potentially targeting vulnerable teens to advertise beauty products. The platform used metrics like if a user recently deleted photos of themselves or made posts referencing feeling insecure or depressed. It would then use this information to advertise products to these vulnerable users. "She said the company was letting advertisers know when the teens were depressed so they could be served an ad at the best time. As an example, she suggested that if a teen girl deleted a selfie, advertisers might see that as a good time to sell her a beauty product as she may not be feeling great about her appearance. They also targeted teens with ads for weight loss when young girls had concerns around body confidence, Wynn-Williams said."\cite{Bellens2025_MetaAdsTeens} This shows that platforms are willing to target vulnerable users to sell a "solution" to their problems.

\section*{Conclusion}
Overall, social media platforms need algorithms to function. Without a system in place to sort out uninteresting and unimportant posts, these platforms would be nearly unusable. These platforms require engagement from their users in order to profit and keep active. However, there is an ethical responsibility to protect users from harmful content. This responsibility is often overlooked or a low priority for companies mainly focused on profit.

With millions of accounts posting daily, it is impossible for moderation to be perfect. For many of these companies, profit is the primary concern and user safety is often overlooked. Because of this, the best way to avoid harmful content is to understand the types of content that may be harmful to you and do your best to avoid them. Information presented in social media posts should not be trusted until verified by a reputable source.

Social media algorithms are not going away. However, some steps may be taken to reduce the harm they cause. Many platforms suffer from a lack of quality moderation. Moderators are often overworked and are unable to deal with the volume of posts on the sites they moderate. AI moderation has become more popular recently, but is often flawed and exploitable. Increasing moderation would reduce the amount of extreme posts that make it past filters. Adjusting algorithms to prioritize positive posts over pure engagement is another way that the spread of harmful content could be limited. Overall, more transparency with how these algorithms work and what gets prioritized would allow users to make informed decisions and understand why they are seeing certain posts. 

\vspace{4\baselineskip}\vspace{-\parskip}
\footnotesize
\bibliographystyle{acm} 
\bibliography{paper}

\end{document}

